% Draft for morletdb.tex, proposed between "Storage, recall and latency" and "Related work".
% Draft 2, 2026-11-11. Author: Semir Maslo (Sulzer-BDiM).
% Public data only (Bosch CNC Machining, CC BY 4.0). Every number comes from the repository:
%   frozen study IU1:      evaluation/industrial_iu1/results/summary.json
%   addenda A1, A2 (post hoc): evaluation/industrial_iu1/results/drift_summary.json, leadin_summary.json

\section{Industrial use}\label{sec:industrial}
Condition monitoring motivated \M{}, but the retrieval results above are seismic. This section reports a first
study on public machine data.

\textbf{Study.} We used the Bosch CNC milling data~\cite{bosch-cnc} (CC BY 4.0): 1,702 runs of 15 operations on three
machines, labelled good or bad per run, tri-axial acceleration in integer counts at 2\,kHz, 178.6\,M samples per axis.
The data give the half-year of a run, not its time. Each run and axis is one recording; envelopes use the product
defaults (windows of 4,096 samples, stride 2,048, $L$ from \texttt{max\_lag\_fraction} 0.075). Machine M01 served to
check the harness; M02 and M03 (1,183 runs, 36 bad) are the confirmation data. The protocol, with every decision rule,
was frozen before loading. Up to eight query windows per run and axis; hits from the query's own run are removed. We
compare bounded lag (BL) with the same scorer at zero lag (ZL) and with nearest windows by RMS, and decide by a 95\%
run bootstrap.

\textbf{Scoring runs against a good history.} Scoring every run by its distance to the nearest window of good runs of
the same machine, operation and axis gave an AUROC for bad against good of 0.925 (95\% interval 0.878--0.967) under BL,
0.875 under ZL and 0.852 under RMS. BL was better than ZL; against RMS the difference was not significant. Asked for
the same operation on another machine, BL returned it in 0.503 of the top ten against 0.412 for ZL (chance 0.076). BL
and ZL shared only 38\% of their hits.

\textbf{Is it drift?} Twenty of the 36 bad runs come from one machine and half-year, so after seeing the results we
checked, with rules written down first, whether the score measures drift between half-years. It does not: with the
history restricted to the query's own half-year the AUROC was 0.933 (0.895--0.965), and over pairs of a bad and a good
run from the same machine and half-year 0.905 (0.853--0.950) under BL against 0.790 under ZL (difference $+0.116$,
0.047--0.184). Without that machine and half-year it was 0.854 (0.768--0.932), and without each run's first two seconds 0.900
(0.836--0.954), so the lead-in does not carry it either. Three restrictions follow. A history
from earlier half-years only did not separate the classes (0.603, 0.419--0.780, 13 bad runs): the signal changes
between half-years by more than a bad run differs from a good one, so the history must be recent, and the data do not
allow a causal test at finer resolution. Run length alone separated better than any signal score (0.982 against 0.925;
inside a machine and half-year 0.988 against 0.905): in these data a bad run is mostly a run of unusual length.
And retrieval of other bad runs was inflated by the half-year: inside the query's own half-year the top ten held bad
runs in 0.254 of positions under BL, 0.212 under ZL and 0.377 under RMS, at a chance level of 0.183 (the frozen
protocol had reported 0.177 against a chance level of 0.026).

\textbf{Time warping.} We stretched query intervals by 1, 2, 5 and 10\% to imitate a feed override and searched for
the source interval. It stayed in the BL top ten for 0.97, 0.76, 0.51 and 0.20 of 100 queries, against 1.00, 0.79,
0.20 and 0.07 for ZL. Bounded lag absorbs small offsets, not changes of speed.

\textbf{Probes.} On an unfiltered gallery of 85,042 windows, probe search returned the exact top ten for all 100
queries, at a p50 of 56\,ms against 141\,ms for exact search and a p95 of 126 against 176\,ms. The frozen rule asked
for half the exact p95, so the speed claim does not hold at this gallery size; it agrees with Section~\ref{sec:limits}
that exact search suits galleries below a few hundred thousand windows.

\textbf{What follows for industrial use.} Lag tolerance improves retrieval on machining vibration as it does on
seismic envelopes: BL beat ZL in every comparison, inside and across half-years. A nearest-neighbour score against a
recent good history separates bad runs without training, but simpler values do as well or better here: RMS for
retrieval and run length for scoring. The envelope is normalised and so blind to level and duration; a product should
offer both next to the shape.
A changed speed needs alignment or segmentation by NC block before search; at 5\% stretch half of the sources are
lost. The study covers one public dataset, run-level labels without time stamps, single axes and synthetic stretching
on one laptop; it makes no claim on diagnosis or on real-time decisions at the machine.

% --- Ledger rows (Claims and evidence) ---
% history score AUROC 0.925 (BL), 0.875 (ZL), 0.852 (RMS); BL vs RMS n.s. & IU1; E, C & run labels, all half-years as history \\
% inside machine and half-year: BL 0.905 vs ZL 0.790; without M02 Feb 2019: 0.854 & IU1-A1; E, C & post hoc, rules written first \\
% earlier half-years only: BL 0.603 (n.s.), 13 bad runs & IU1-A1; E, C & half-year resolution only \\
% without each run's first 2 s: BL 0.900 & IU1-A2; E, C & query side only, 2 s assumed \\
% run length alone: AUROC 0.982, better than BL & IU1-A1; E, C & run length from file length \\
% other bad runs in top 10 inside the half-year: BL 0.254, ZL 0.212, RMS 0.377, chance 0.183 & IU1-A1; E, C & replaces 0.177 vs 0.026 of IU1 \\
% same operation on other machine 0.503 vs 0.412 & IU1; E, C & axis filter only \\
% stretch 1/2/5/10\%: source in top 10 0.97/0.76/0.51/0.20 & IU1; E, C & synthetic stretch \\
% probes recall 1.0, p95 126 vs 176 ms at 85k windows & IU1; P, C & laptop, Docker \\

% --- Bibliography entry ---
% @article{bosch-cnc, author={Tnani, Mohamed-Ali and Feil, Michael and Diepold, Klaus},
%   title={Smart Data Collection System for Brownfield {CNC} Milling Machines: A New Benchmark Dataset for
%   Data-Driven Machine Monitoring}, journal={Procedia CIRP}, volume={107}, pages={131--136}, year={2022},
%   doi={10.1016/j.procir.2022.04.022}}
